% =============================================================================
% Project : THE_PHILOSOPHY_OF_VIABILITY
% Book    : Book I — Ontology of Viability
% Book ID : VS.B01
% Part    : Part I — Identity and Manifestation (VS.B01.P01)
% Chapter : Chapter 08 — Comparative Philosophical Examination
% Chapter ID : VS.B01.P01.C08
%
% File    : 4_ASSETS/diagrams/VS.B01.P01.C08.S02.SS03.P006.DG01.parmenidean_routes_and_deduction.tex
%
% Section / Chapter / Asset:
%   Diagram DG01 — Parmenidean Routes and Deduction
%
% Structural Identifier : VS.B01.P01.C08.S02.SS03.P006.DG01
% Asset Type : DG
%
% Purpose :
%   To distinguish the two ways introduced in B2 from the mortal route in B6,
%   and to represent the inferential movement from the affirmation of what-is
%   to the principal signs in B8.
%
% Representational limit :
%   The solid arrows mark the argumentative movement; the dashed arrows mark
%   routes that the poem presents as unavailable or unstable. The
%   reconstruction does not settle the interpretive disputes concerning the
%   fragments.
%
% Dependencies :
%   TikZ; libraries used: arrows.meta, calc, positioning, fit. All are loaded
%   by the project preamble; no per-file package or library loading.
%
% =============================================================================

\begin{figure}[htbp]
    \centering
    \hyphenpenalty=10000 \exhyphenpenalty=10000\relax  % no hyphenation inside diagram labels

    \begin{tikzpicture}[
        font=\small,
        node distance=0.8cm,
        % --- base box ---
        box/.style={
            draw, rounded corners=3pt, align=center,
            text width=3.6cm, minimum height=1.5cm,
            inner sep=5pt, line width=0.7pt
        },
        % --- node roles ---
        route/.style={box, fill=blue!8, draw=blue!55!black},
        deduction/.style={box, fill=blue!16, draw=blue!55!black,
                          minimum height=1.3cm},
        sign/.style={box, fill=green!8, draw=green!40!black,
                     minimum height=1.0cm},
        excluded/.style={box, dashed, fill=gray!10, draw=black!60,
                         minimum height=1.3cm},
        % --- frames and labels ---
        vsgroup/.style={draw=black!45, dashed, rounded corners=5pt,
                        inner sep=8pt, line width=0.5pt},
        grouplabel/.style={font=\footnotesize\itshape, text=black!75,
                           align=center},
        % --- arrows ---
        directed/.style={-{Stealth[length=2.5mm]}, line width=0.9pt},
        optional/.style={-{Stealth[length=2.5mm]}, line width=0.9pt, dashed}
    ]

    % ---------------------------------------------------------------
    % Ways of inquiry (B2) and the mortal route (B6)
    % ---------------------------------------------------------------
    \node[route] (being) {
        \textbf{B2: What-is}\\[2pt]
        It is and cannot not be
    };
    \node[route, right=0.8cm of being] (nonbeing) {
        \textbf{B2: What-is-not}\\[2pt]
        It is not and must not be
    };
    \node[route, right=0.8cm of nonbeing] (mortal) {
        \textbf{B6: Mortal route}\\[2pt]
        Confused treatment of being and not-being
    };
    \node[vsgroup, fit=(being)(nonbeing)(mortal)] (routes) {};
    \node[grouplabel, above=0.1cm of routes] {Ways of inquiry};

    % ---------------------------------------------------------------
    % Consequences of each route
    % ---------------------------------------------------------------
    \node[deduction, below=1.3cm of being] (b8) {
        \textbf{B8: Deduction}\\[2pt]
        What-is must be
    };
    \node[excluded, below=1.3cm of nonbeing] (unavailable) {
        \textbf{Not available}\\[2pt]
        No attainable object of inquiry
    };
    \node[excluded, below=1.3cm of mortal] (wandering) {
        \textbf{Wandering inquiry}\\[2pt]
        Unstable judgment
    };

    % ---------------------------------------------------------------
    % Principal signs of what-is (B8)
    % ---------------------------------------------------------------
    \node[sign, below=1.3cm of b8] (s1) {Ungenerated and imperishable};
    \node[sign, below=0.3cm of s1] (s2) {Whole and continuous};
    \node[sign, below=0.3cm of s2] (s3) {Unmoved and complete};
    \node[vsgroup, fit=(s1)(s2)(s3)] (signs) {};
    \node[grouplabel, right=0.2cm of signs, anchor=west, text width=2.6cm,
          align=left] {Principal signs of what-is (B8)};

    % ---------------------------------------------------------------
    % Relations
    % ---------------------------------------------------------------
    \draw[directed] (being)   -- (b8);
    \draw[directed] (b8)      -- (signs);
    \draw[optional] (nonbeing) -- (unavailable);
    \draw[optional] (mortal)   -- (wandering);

    % ---------------------------------------------------------------
    % Legend and interpretive caution
    % ---------------------------------------------------------------
    \coordinate (axis) at (nonbeing |- signs.south);
    \node[below=0.5cm of axis, font=\scriptsize] (legend) {
        \tikz[baseline=-0.5ex]\draw[directed] (0,0) -- (0.7,0);~argumentative movement
        \quad\quad
        \tikz[baseline=-0.5ex]\draw[optional] (0,0) -- (0.7,0);~route presented as unavailable or unstable
    };
    \node[below=0.35cm of legend, text width=11cm, align=center,
          font=\footnotesize\itshape] {
        The relation between the mortal route, the deductions concerning
        what-is, and the cosmology of Doxa remains a matter of
        interpretation. The mortal route is not equated with absolute
        non-being.
    };

    \end{tikzpicture}

    \caption{The ways of inquiry in B2 and B6 and the deductive structure
    of B8. The solid arrows mark the argumentative movement; the
    dashed arrows mark routes that the poem presents as
    unavailable or unstable. The reconstruction does not settle
    the interpretive disputes concerning the fragments.}
    \label{fig:ch08-parmenides-routes-deduction}
\end{figure}

% =============================================================================
% END OF DIAGRAM DG01
% =============================================================================
